% =====================================================
% 题型：三角形面积、海伦公式与四边形面积最值解答题 (q_lyb_01_19_heron_quadrilateral.tex)
% =====================================================
\newcommand{\qLYBZeroOneNineteenStem}{在 \(\triangle ABC\) 中，内角 \(A, B, C\) 的对边分别为 \(a, b, c\)。\\
(1) 若 \(a = 1\), \(b = 2\), \(c = \sqrt{7}\)，求 \(\triangle ABC\) 的面积;\\
(2) 设 \(p = \dfrac{1}{2}(a + b + c)\)，证明：\(\triangle ABC\) 的面积 \(S = \sqrt{p(p-a)(p-b)(p-c)}\);\\
(3) 若平面凸四边形 \(ABCD\) 中，\(AB = BC = AD = 1\), \(CD = 2\)，求四边形 \(ABCD\) 面积的最大值。}

\newcommand{\qLYBZeroOneNineteenAnswer}{
(1) \(\dfrac{\sqrt{3}}{2}\); (2) 见解析; (3) \(\dfrac{3\sqrt{3}}{4}\).
}

\newcommand{\qLYBZeroOneNineteenSolution}{%
  (1) 解：
  在 \(\triangle ABC\) 中，由余弦定理可得：
  \[ \cos C = \frac{a^2 + b^2 - c^2}{2ab} = \frac{1^2 + 2^2 - (\sqrt{7})^2}{2 \cdot 1 \cdot 2} = \frac{1 + 4 - 7}{4} = -\frac{1}{2} \]
  因为 \(C \in (0, \pi)\)，所以：
  \[ \sin C = \sqrt{1 - \cos^2 C} = \sqrt{1 - \left(-\frac{1}{2}\right)^2} = \frac{\sqrt{3}}{2} \]
  因此，\(\triangle ABC\) 的面积为：
  \[ S = \frac{1}{2}ab\sin C = \frac{1}{2} \cdot 1 \cdot 2 \cdot \frac{\sqrt{3}}{2} = \frac{\sqrt{3}}{2} \]
  故 \(\triangle ABC\) 的面积为 \(\dfrac{\sqrt{3}}{2}\)。\\
  
  (2) 证明：
  三角形的面积公式为：
  \[ S = \frac{1}{2}ab\sin C \implies S^2 = \frac{1}{4}a^2b^2\sin^2 C = \frac{1}{4}a^2b^2(1 - \cos^2 C) \]
  由余弦定理 \(\cos C = \dfrac{a^2 + b^2 - c^2}{2ab}\) 代入上式：
  \[ S^2 = \frac{1}{4}a^2b^2 \left[ 1 - \left(\frac{a^2 + b^2 - c^2}{2ab}\right)^2 \right] = \frac{1}{4}a^2b^2 \cdot \frac{4a^2b^2 - (a^2+b^2-c^2)^2}{4a^2b^2} \]
  \[ S^2 = \frac{1}{16} \left[ (2ab)^2 - (a^2+b^2-c^2)^2 \right] \]
  利用平方差公式展开：
  \[ S^2 = \frac{1}{16} (2ab + a^2 + b^2 - c^2)(2ab - a^2 - b^2 + c^2) \]
  \[ S^2 = \frac{1}{16} \left[ (a+b)^2 - c^2 \right] \left[ c^2 - (a-b)^2 \right] \]
  再次利用平方差公式：
  \[ S^2 = \frac{1}{16} (a+b+c)(a+b-c)(c+a-b)(c-a+b) \]
  因为 \(p = \dfrac{a+b+c}{2}\)，所以：
  \(a+b+c = 2p\), \\
  \(a+b-c = (a+b+c) - 2c = 2p - 2c = 2(p-c)\), \\
  \(c+a-b = (a+b+c) - 2b = 2p - 2b = 2(p-b)\), \\
  \(c-a+b = (a+b+c) - 2a = 2p - 2a = 2(p-a)\)。\\
  代入上式，得：
  \[ S^2 = \frac{1}{16} \cdot (2p) \cdot 2(p-c) \cdot 2(p-b) \cdot 2(p-a) = p(p-a)(p-b)(p-c) \]
  两边开平方根（由于面积 \(S > 0\)）：
  \[ S = \sqrt{p(p-a)(p-b)(p-c)} \]
  海伦公式得证。\\
  
  (3) 解：
  连接对角线 \(BD\)。
  在凸四边形 \(ABCD\) 中，设 \(\angle BAD = A\)，\(\angle BCD = C\)。
  因为 \(AB = BC = AD = 1\), \(CD = 2\)，
  在 \(\triangle ABD\) 中，由余弦定理得：
  \[ BD^2 = AB^2 + AD^2 - 2AB \cdot AD \cos A = 1^2 + 1^2 - 2 \cdot 1 \cdot 1 \cos A = 2 - 2\cos A \]
  在 \(\triangle BCD\) 中，由余弦定理得：
  \[ BD^2 = BC^2 + CD^2 - 2BC \cdot CD \cos C = 1^2 + 2^2 - 2 \cdot 1 \cdot 2 \cos C = 5 - 4\cos C \]
  由于对角线 \(BD\) 是公共边，所以：
  \[ 2 - 2\cos A = 5 - 4\cos C \implies \cos A = 2\cos C - \frac{3}{2} \]
  四边形 \(ABCD\) 的面积 \(S_{ABCD}\) 为 \(\triangle ABD\) 与 \(\triangle BCD\) 面积之和：
  \[ S_{ABCD} = S_{\triangle ABD} + S_{\triangle BCD} = \frac{1}{2}AB \cdot AD \sin A + \frac{1}{2}BC \cdot CD \sin C = \frac{1}{2}\sin A + \sin C \]
  我们要最大化 \(S_{ABCD} = \frac{1}{2}\sin A + \sin C\)。
  因为 \(ABCD\) 是平面凸四边形，所以 \(A, C \in (0, \pi)\)，即 \(\sin A > 0, \sin C > 0\)。
  为了利用已有的约束条件 \(\cos A - 2\cos C = -\dfrac{3}{2}\) 且避开求导，我们将面积表达式乘以 2，设为：
  \[ 2S_{ABCD} = \sin A + 2\sin C \]
  根据三角恒等式，我们将约束方程和目标方程两边平方并相加，利用平方和为 1 以及和角公式来化简：
  \[ (\cos A - 2\cos C)^2 + (\sin A + 2\sin C)^2 = \left(-\frac{3}{2}\right)^2 + (2S_{ABCD})^2 \]
  将等号左边展开：
  \[ \text{LHS} = (\cos^2 A + \sin^2 A) + 4(\cos^2 C + \sin^2 C) - 4(\cos A\cos C - \sin A\sin C) \]
  代入 \(\cos^2 x + \sin^2 x = 1\) 以及两角和的余弦公式 \(\cos(A+C) = \cos A\cos C - \sin A\sin C\)，得：
  \[ \text{LHS} = 1 + 4 - 4\cos(A+C) = 5 - 4\cos(A+C) \]
  由此可得：
  \[ \frac{9}{4} + 4S_{ABCD}^2 = 5 - 4\cos(A+C) \implies 4S_{ABCD}^2 = \frac{11}{4} - 4\cos(A+C) \]
  因为余弦函数 \(\cos(A+C) \ge -1\)，所以有：
  \[ 4S_{ABCD}^2 \le \frac{11}{4} - 4(-1) = \frac{27}{4} \]
  由于面积 \(S_{ABCD} > 0\)，可得：
  \[ S_{ABCD}^2 \le \frac{27}{16} \implies S_{ABCD} \le \frac{3\sqrt{3}}{4} \]
  当且仅当 \(\cos(A+C) = -1\) 时等号成立。
  由于 \(A, C \in (0, \pi) \implies A+C \in (0, 2\pi)\)，此时必有 \(A + C = \pi\)（即四边形 \(ABCD\) 为圆内接四边形）。
  将 \(A = \pi - C\) 代入约束条件 \(\cos A = 2\cos C - 1.5\)：
  \[ \cos(\pi - C) = 2\cos C - 1.5 \implies -\cos C = 2\cos C - 1.5 \implies 3\cos C = 1.5 \implies \cos C = \frac{1}{2} \]
  因为 \(C \in (0, \pi)\)，所以 \(C = \dfrac{\pi}{3}\)，此时 \(A = \pi - C = \dfrac{2\pi}{3}\)。
  由于 \(A, C \in (0, \pi)\)，此状态完全符合题意。
  因此，四边形 \(ABCD\) 面积的最大值为 \(\dfrac{3\sqrt{3}}{4}\)。%
}
